\section{Алфавит, слово и язык. Кодирование и его свойства. Методы описания множества сообщений. Алфавитное кодирование.}

\subsection*{Алфавит, слово и язык}

\paragraph{Определения}
\begin{itemize}
    \item \textbf{Алфавит} ($A$) — конечное множество символов (букв).
    \item \textbf{Слово} — любая конечная последовательность символов алфавита $A$.
    \item \textbf{Пустое слово} ($\varepsilon$) — слово, не содержащее ни одного символа. Длина $| \varepsilon | = 0$.
    \item \textbf{Длина слова} $|w|$ — количество символов в слове.
    \item \textbf{Множество слов:}
    \begin{itemize}
        \item $A^*$ — множество всех слов в алфавите $A$, включая пустое слово.
        \item $A^+$ — множество всех непустых слов ($A^* \setminus \{ \varepsilon \}$).
    \end{itemize}
    \item \textbf{Язык} ($L$) — любое подмножество множества всех слов ($L \subseteq A^*$).
\end{itemize}

\subsection*{Методы описания множества сообщений}
Сообщения — это слова некоторого языка. Описать язык можно тремя способами:
1. \textbf{Перечисление:} (только для конечных языков) $\{ \text{слово}_1, \text{слово}_2, \dots \}$.
2. \textbf{Порождающие процедуры:} использование грамматик (правил вывода слов).
3. \textbf{Распознающие процедуры:} использование автоматов или регулярных выражений, которые отвечают «да/нет» на вопрос, принадлежит ли слово языку.

\subsection*{Кодирование и его свойства}

\paragraph{Кодирование} — это отображение $f$ элементов одного множества (исходных символов $S$) в слова другого алфавита $B$ (кодовые слова).
$$f: S \to B^*$$

\paragraph{Свойства кодирования:}
\begin{enumerate}
    \item \textbf{Инъективность:} разным исходным символам соответствуют разные кодовые слова.
    \item \textbf{Однозначная декодируемость:} любое слово в кодовом алфавите (последовательность кодовых слов) может быть разложено на исходные символы единственным способом.
    \item \textbf{Префиксное свойство (Условие Фано):} никакое кодовое слово не является началом (префиксом) другого кодового слова.
    \item \textbf{Оптимальность:} минимизация средней длины кодового слова (например, код Хаффмана).
\end{enumerate}

\subsection*{Алфавитное кодирование}

\paragraph{Схема кодирования}
Пусть задан исходный алфавит $S = \{s_1, s_2, \dots, s_n\}$ и целевой алфавит $B$. \textbf{Схема алфавитного кодирования} — это таблица соответствия:
\begin{itemize}
    \item $s_1 \to w_1$
    \item $s_2 \to w_2$
    \item \dots
    \item $s_n \to w_n$
\end{itemize}
где $w_i \in B^*$. Кодирование сообщения $s_{i_1}s_{i_2}\dots s_{i_k}$ осуществляется простой конкатенацией кодовых слов: $w_{i_1}w_{i_2}\dots w_{i_k}$.

\paragraph{Виды кодов:}
\begin{itemize}
    \item \textbf{Равномерные коды:} все кодовые слова $w_i$ имеют одинаковую длину (пример: код ASCII, где на символ 8 бит).
    \item \textbf{Неравномерные коды:} длины кодовых слов различаются (пример: азбука Морзе, код Хаффмана).
\end{itemize}

\subsection*{Примеры}
\begin{enumerate}
    \item \textbf{Двоичное кодирование:} $A = \{a, b, c, d\}$, схема: $a \to 00, b \to 01, c \to 10, d \to 11$. Это равномерный префиксный код.
    \item \textbf{Код Хаффмана:} для часто встречающихся символов назначаются короткие коды, для редких — длинные. Код является префиксным по построению.
    \item \textbf{Азбука Морзе:} $s_1 \to \cdot, s_2 \to -$. Не является префиксным кодом в чистом виде (требует разделителя — паузы между буквами), поэтому строго говоря, алфавитом кодирования в Морзе является $\{\cdot, -, \text{пауза}\}$.
\end{enumerate}